\documentclass{article}
\usepackage{cleveref}
\usepackage{tabularray}
\usepackage{graphicx}
\usepackage{amsmath}

\begin{document}

% === L10: 标题未 Title Case ===
\section{research background and motivation}
\label{background}

% === L09: 缩写首次出现未展开 ===
FHE is a powerful cryptographic technique.
NTT is the key computation kernel.

% === L01: 用了 \ref 而不是 \cref ===
As shown in Figure \ref{fig:overview},
the architecture is illustrated.

% === L02/L03: 缺 ~ 非断行空格 ===
See Table \cref{tab:results} for details.
As discussed in \cite{cinnamon2025}.

% === L04/L05/L06: 浮动体综合问题 ===
\begin{figure}
\centering
\label{overview}
\caption{System Overview}
\end{figure}

% === L07: 用了 tabular 而不是 tblr ===
\begin{table}[htbp]
\caption{Performance Comparison}
\label{tab:results}
\begin{tabular}{lcc}
\hline
Method & Speedup & Area \\
\hline
Ours & 3.2x & 1.5mm2 \\
\hline
\end{tabular}
\end{table}

% === L08: 未引用的 label ===
\begin{equation}
E = mc^2
\label{eq:unused}
\end{equation}

% === L10: 另一个标题 ===
\subsection{our proposed method}

% === L11: 缩写后句间距 ===
In Fig. 1, we show the results.
Et al. proposed a method.

% === L13: 连续多空格 ===
This is  a  test  sentence.
Another sentence with trailing spaces.    

\bibliography{lint_buggy}
\end{document}
